\documentclass[11pt]{article}
\usepackage{graphicx} % Required for inserting images
\usepackage{hyperref} % Required for hyperlinks

% Beautiful theorems, definitions etc.
\usepackage[most]{tcolorbox}
\usepackage{amsmath, amsthm}

% This command forces dotted lines for sections in the article class
\usepackage{tocloft}
\renewcommand{\cftsecleader}{\cftdotfill{\cftdotsep}}

\title{Massachusetts Institute of Technology - Introduction to Algorithms (6.1210)}
\author{Xavier Henry}
\date{September 2026}

% --- 1. Page Setup & Geometry ---
\usepackage[utf8]{inputenc}
\usepackage[margin=1in]{geometry} % Fixes the massive default LaTeX margins

% --- 2. Math & Science Essentials ---
\usepackage{amsmath, amssymb, amsthm} % Essential for formulas, symbols, and proofs
\usepackage{bm} % For bold math symbols
\usepackage{enumitem} % for lists

\usepackage[normalem]{ulem} % for strikethrough!

% --- 3. Beautiful Callout Boxes (Perfect for Notes!) ---
% 3.1. Define a standard theorem counter that resets with the section
\newtheorem{thm}{Theorem}[section]
\newtheorem{defnt}{Definition}[section]
\newtheorem{exmpl}{Example}[section]
\newtheorem{corol}{Corollary}[section]

% 3.2. Create the color boxes and link them to those counters using "use counter"
\newtcolorbox[use counter=defnt]{definitionbox}[1]{
    colback=blue!5!white, 
    colframe=blue!75!black, 
    fonttitle=\bfseries, 
    title=Definition \thedefnt: #1, % Automatically prints "Definition X.Y: Your Title"
    arc=4mm
}

\newtcolorbox[use counter=thm]{theorembox}[1]{
    colback=purple!5!white,    % 5% purple background (soft pastel)
    colframe=purple!75!black,  % Deep purple border and title bar 
    fonttitle=\bfseries, 
    % This line checks if you gave a title. If not, it skips the colon
    title={Theorem \thethm\ifstrempty{#1}{}{ : #1}},
    arc=4mm
}

\newtcolorbox[use counter=exmpl]{examplebox}[1]{
    colback=green!3!white,    % 5% purple background (soft pastel)
    colframe=green!80!black,  % Deep purple border and title bar 
    fonttitle=\bfseries, 
    % This line checks if you gave a title. If not, it skips the colon
    title={Example \theexmpl\ifstrempty{#1}{}{ : #1}},
    arc=4mm
}

\newtcolorbox[use counter=corol]{corollarybox}[1]{
    colback=yellow!5!white,    % 5% purple background (soft pastel)
    colframe=yellow!75!black,  % Deep purple border and title bar 
    fonttitle=\bfseries, 
    % This line checks if you gave a title. If not, it skips the colon
    title={Corollary \thecorol\ifstrempty{#1}{}{ : #1}},
    arc=4mm
}

\begin{document}

\maketitle

\tableofcontents \newpage

\section{Lecture 1}
\begin{definitionbox}{Algorithm}
    An algorithm is a step by step procedure for solving a computational, problem. A computational problem takes inputs, and for each input, outputs one or more resulting outputs.
\end{definitionbox}

\subsection{Integer Addition}

Consider the preschool addition algorithm $x + y$ where we add 1 to $x$, $y$ times. The running time of this is $\Theta(10^n)$. Prove this! \\
Our grade school algorithm, wherein we line up the digits and add (assuming we've memorized the single-digit addition table), takes linear time, $\Theta(n)$

\begin{theorembox}{Integer Addition Lower Bound}
    Every correct algorithm for integer addition requires $\Omega(n)$ time.
\end{theorembox}

\begin{proof}
    Our output must have at least $n$ digits.
\end{proof}

Alternatively, consider the following (more subtle) proof
\begin{proof}
    For our algorithm to be correct, we must read every input digit, of which there are $\Omega(n)$.
    Suppose we have a magical correct addition algorithm $f$ that does not read every input digit. Without loss of generality, suppose the digit that is not touched is $x_0$. There exists  $x,y$ and $x',y'$ which differ on only $x_0$ st. $x+y \neq x'+y'$ but $f(x,y) = f(x',y')$.
\end{proof}

\subsection{Integer Multiplication}
We now consider the preschool integer multiplication, where we add $x$ to itself $y$ times. This is exponential time. \\ % WHY???
Consider also the grade school multiplication algorithm. This is the standard algirithm of lining up the digits and multiplying them with our multiplication tables. We know that each digit in $x$ must touch each digit in $y$, so this must be $\Theta(n)$. \\
Importantly, is there a better algorithm that is correct, yet faster? \\
We know the size of our output is linear $O(n)$ and the size of our input is $O(n)$.

\subsubsection{Divide and Conquer}
A divide and conquer algorithm divides our problem into subproblems, recursively solves each subproblem, and assebles into a solution.
We can use such an algorithm for integer multiplication (figure this out)!

\begin{align*}
    A &= mult(x_l, y_l) \\
    B &= mult(x_h, y_l) \\
    C &= mult(x_l, y_h) \\
    D &= mult(x_h, y_h) \\
    T(n) &= \text{worst-case time on n-digit inputs} \\
    T(n) &\leq 4T(n/2)+Cn \\
    T(n) &\leq O(n^2) % go back and solve this
\end{align*}

\subsubsection{Karatsuba's Algorithm}
Instead, let's do:
\begin{align*}
    A &= mult(x_l, y_l) \\
    D &= mult(x_h, y_h) \\
    E &= mult(x_l + x_h, y_l + y_h) \\
\end{align*}

And our output: $A + 10^{n/2}(E-A-D) + 10^nD$. This becomes clear if we FOIL.

Now we prove our running time.
\begin{proof}
    \begin{align*}
    T(n) &\leq 3T(n/2 + 1) + Cn \\
    T(n) &\leq n^{\log_2{3}} \\
         &= O(n^{1.6})
    \end{align*}
\end{proof}
With some analysis of the algorithm, we can prove that it is tight. \\
It's important to note, that though Karatsuba is better asymptotically, for smaller input values, the speed vs. grade school multiplication depends heavily on our implementation of the algorithm.

\section{Lecture 2: Recursion}

\subsection{Sorting}
Input: Array $A$ \\
Output: Permutation $\sigma$ such that $A \circ \sigma$ is monotone.

\subsubsection{Selection Sort}
We will analyze a recursive implementation of selection sort.
\begin{proof}
    This is just the implementation.
    Base case: If $|A| \leq 1$ return $A$ \\
    Otherwise: Let $m :=$ Arg ,min $A$. Swap $A[0],A[m]$ \\
    Call selection sort on $A[1:]$
\end{proof}
You can see this is very similar to induction. We now write a proof of correctness
\paragraph{Proof of Correctness}
Correctness: Inductive Preidcate $P(n)$, Alg correct on size $n$ input.
\begin{proof}
    Base Case: If $|A| \leq 1$ already sorted, so return $A$ (correct result). \\
    Inductive Step: Suppose $|A| > 1$ and $P(n-1)$. We want to prove that we end up with a permutation of $A$ we call $A'$ (1) and that this permutation is monotone (2). To prove (1), we ensure that at each step, we are applying a permutation. Finding the min of $A$ has no write, performing a swap is a transposition and $SS(A[1:])$ permutes by our inductive hypothesis, since we assume $P$ for smalller values. Therefore, $A'$ is a permutation of $A$. (proved 1) \\
    We swap the min of $A$ to $A'[0]$. By our inductive hypothesis, $A'[1:]$ is monotone. $A'[1]$ is in $A$ so is no smaller than min $A$. \\
    $A'[0] \leq A'[1] \leq ...$ so $A'$ is monotone (proved 2) \\
    By induction, our implementation of selection sort is correct.
\end{proof}
\paragraph{Efficiency and Proving Bounds}
Quantify input
\begin{definitionbox}{Worst-Case Runtime}
    The worst-case runtime $T(n)$ of $A$ is our maximum over all $x$ of size $n$ of number of opperations used by $A(x)$. $A$ has worst case runtime $O(f)$ means for sufficiently large $x$, $A(x)$ uses $\leq$ constant factor for more than $f(|x|)$ operations. Finding this worst-case runtime gives us a lower bound.\\
    If we want to prove an upper bound on our worst-case runtime, we want to prove that $A(x)$ always takes $O(f(|x|))$ time. As said before, to prove a lower bound we need \textbf{one} input of each length for which $A(x)$ is slow. If we choose a particularly fast input for $A$ to lower bound our algorithm, this is wrong.
\end{definitionbox}
\begin{definitionbox}{Model of Computation and Atomic Operations}
    A model of computation is an abstraction of a computer. We use what is called the word-RAM model, which closely mimicks the operations of an actual computer. We use $w$ bit words as a unit of memory. Our atomic operations are $+,-,*,/, etc.$, AND, OR, NOT, comparisions, $...$. Any bianry operation is atomic. We can read and write from/to specific memory addresses. We also have control flow, so things like IF, WHILE, RETUR. We can do all these atomic operations in constant time.
\end{definitionbox}
To analyze the runtime of selection sort, we need a recurrence relation.
\begin{align*}
    T(n) \in \Theta(1) \text{if} n \leq 1 \\
    T(n) &\leq \Theta(n) + \Theta(1) + T(n-1) \\ 
    T(n) &\leq O(n^2)
\end{align*}
We ignore the largest half of elements $n/2$ times. \\
Lower bound: $n/2$ elements read $n/2$ times each $\Omega(n^2)$. Say we have a monotonically decreasing list that we want to make monotonically increasing. By the time selection sort reaches $\lfloor n/2 \rfloor$, we have sorted the first half of the list, and swapping has sorted the second half of the list. This proves our $n^2$ lower bound, since we've done at least $n/2$ work for $n/2$ elements.
\subsubsection{Peak Finding}
\begin{definitionbox}{Peak}
    A peak of array $A$ is an index $i$ such that $A[i]$ is at least as large as its neighbours (local maximum)
\end{definitionbox}

\paragraph{Peak-finding Problem}
Input: Array $A$ (non-empty) \\
Output: Peak of $A$ (any) \\
Edge cases of input should be considered carefully in the problem definition!

\paragraph{Algorithm 1: Exhaustive Search} As counter intuitive as this may seem, this is a perfectlty valid algorithm definition. \\
We first prove correctness.
\begin{proof}
    Correctness by exhaustion.
\end{proof}
Runtime:
\begin{proof}
    Upper: By definition we do $\Theta(1)$ check at most $O(n)$. \\
    Lower: Suppose $A$ strictly monotone. Our check fails $\Omega(n)$ times. (notice here we picked a super slow input and used that as our lower bound!)
\end{proof}
\paragraph{Algorithm 2: Binary Search}
Recursion: discard half (constant fraction) of input quickly. \\
Idea: check the midpoint (if peak, then return midpoint), go uphill
\begin{proof}
    Here's our implementation: \\
    PF(A): If length 1 outpit 0 \\
    Otherwise, Let $j= \lfloor n/2 \rfloor$. If peak output $j$. If $A[j-1] \geq A[j]$ recurse on $A[:j]$. Otherwise, recurse on $A[j+1:]$.
\end{proof}
Correctness
\begin{proof}
    Base case: if singleton index $0$ is peak vacuously. \\
    Inductive case: Suppose $n > 1$, $PF$ is correct on inputs smaller than $n$. \\
    Case 1: $j$ is peak. Correctly return $j$. \\
    Case 2: Assume without loss of generality $A[j-1] \geq A[j]$. \\
    WTP: returned peak of $A[:j]$ is also a peak in $A$. \\
    Case 1: $k < j-1$ then neighbours are same, so still peak. \\
    Case 2: Case 1 + guard???
\end{proof}
TODO: runtime analysis.

\section{Lecture 3: Divide and Conquer, Iteration}
\begin{itemize}
    \item divide
    \item conquer - solve all of the subproblems recursively
    \item recombine - recombine all the solutions to subproblems together to get our final solution
\end{itemize}
\subsection{Merge Sort}
Divide input into two halves. Recursively sort both halves then marge sorted subarrays into a sorted array. Let's make a problem statement
\begin{proof}
    MERGE \\
    INPUT: two arrays $A,B$ \\
    PRECONDITION: $A+B$ sorted \\
    OUTPUT: an array $C$ which is the sorted concatenation of $A$ and $B$ \\
    Initialize two indices $i,j$. Initialize empty $C$ of size $2n$. \\
    While $i+j < 2n$: \\
    \{ \\
        if $A[i] \leq B[j]$ then $C[i+j] := A[i]$, increment $i$ \\
        otherwise, $C[i+j]:=B[j]$, increment $j$ \\
    \} \\
    Output $C$
\end{proof}
\subsubsection{Proof of Correctness}
We could try to use induction, but this would cause us something called buildup error. That is, with induction it's not clear whether we are actually building toward an answer. Here, it's best to use our \textbf{invariants}! But how do we do this?
\begin{itemize}
    \item Model loop as state machine
    \item Think of a \textbf{state} as a snapshot of memory
    \item Start state: initialization
    \item transitions: executions of loop body
    \item final states: loop \textbf{guard} (What is a loop guard??? -> it's just the while loop condition lol)
\end{itemize}
When doing a proof with invariants, we do not necessarily want to prove correctness directly. Instead, we want to prove partial correctness. That is, if the algorithm terminates, it produces the correct answer. So what do we need for our invariant proof?
\begin{itemize}
    \item Inductive predicate: Choose correct invariant(s)
    \item Initialization: Prove that the invariant(s) are true at initialization (our start state)
    \item Preservation: we need to prove that our invariants are preserved properties throughout the operation of our state machine.
    \item Prove invariants $\implies$ partial correctness
    \item Runtime analysis (proves termination)
\end{itemize}
We first prove partial correctness of our \textbf{merge} algorithm (putting the two lists together).
\begin{proof}
    \textbf{Partial Correctness of Merge} $A,B$ read only \\
    I5: $C[:i+j]$ is the concatenation of $A[:i]+B[:j]$ \\
    I6: $C[:i+j]$ is sorted \\
    I3: $A[i] \geq C[:i+j]$ \\
    I4: $B[j] \geq C[:i+j]$ \\
    I1: $0 \leq i \leq n$ \\
    I2: $0 \leq j \leq n$ \\
    \textbf{Initialization}: All vacuous since $i=j=0$ \\
    \textbf{Preservation}: WLOG $A[i] \leq B[j]$ \\
    I3: WTP $A[i+1] \geq C[:i+j+1]$ \\
    By precondition, $A[i+1 \geq A[i] \geq C[i+j] $ \\
    I4: Because $B[j] \geq A[i]$, as above \\
    I5: Adding $A[i]$ to both prefix of $A$ and of $C$ \\
    I6: follows from I3,I6 \\
    \textbf{Partial Correctness}: Follows from I5, I6, and the loop guard. When our loop guard fails, $C$ will be sorted and will be a permutation of all of $A$ and all of $B$. 
\end{proof}

In order to prove our merge sort algorithm (recurisve), we need to use strong induction. Prove this later
\begin{proof}
    Strong induction.
\end{proof}

Let's do runtime analysis now!
\begin{proof}
    $i+j$ starts at $0$, increments with every loop body $\leq 2n$ by loop guard, so $2n$ iterations $\Theta(1)$ per $\Theta(n)$ total.
    \begin{align*}
        T(n) &\leq 2T(n/2)+O(n) \\
        T(n) &\leq O(nlogn)
    \end{align*}
    and we have our upper bound. Now, for our lower bound. Take any valid input of length $n$. Every element is merged $\Omega(logn)$ times and our merge algorithm takes linear time so $\Omega(nlogn)$.
\end{proof}

\subsection{Inversion Counting}
Inversion: pair $(i,j)$ such that $i < j$ and $A[j]<A[i]$. \\
\textbf{Alg1:} Brute force, Runtime $\Theta(n^2)$, Correctness: exhaustive. \\
But notice that inversion counting is roughly how sorted is $A$? In fact, it's an interesting problem to sort our array, then count the number of inversions. In this case, making our problem a bit harder makes our ideal algorithm much faster. \\
Idea: Suppose $A=A[:m]+A[m:]$ either $i<j<m$ (or $m\leq i \leq j$ which we can count recursively) or $i<m\leq j$ \\
While merging, if we take from right half, add $m-i$, $i$ being our iteration index.

\subsection{2 Dimensional Peak Finding}
TODO: This was taught in one of the (very old) 6.006 lectures! You should totally watch his coverage of this.
Input: $n \times m$ array $A$. \\
Output: any peak in $A$
\paragraph{Exhaustive Search}: $\Theta(nm)$ time (Is this linear?)

\subsubsection{Divide and Conquer}
Look at mid col $j$. Find a peak in $j$, then go uphill. Let's find a peak that is the maximum in its column (this prevents an edge case with our boundary).

\paragraph{Runtime}:
\begin{align*}
    T(m,n) &\leq T(m,n/2)+ \Theta(m) \\
           &\leq O(mlogn) \\
\end{align*}
Take $N:=mn$ then it is $O(\sqrt{N}\log N)$

\section{Lecture 4: Reduction and Randomization}

\subsection{Review}
Divide and conquer; three techniques:
\begin{itemize}
    \item Loop invariants: we start with an initial state, and as long as a loop guard held true,
        we stay in the loop and change things (while maintaining our invariants, which we prove)
        and give an output.
    \item Make the problem harder, so that we move from a specific case of something to a more
        general problem that is easier to analyze and/or solve
    \item Recurrences - \textbf{Master's Theorem!} You should \sout{probably} DEFINITELY memorize 
        this!
\end{itemize}

\subsection{Another Loop: Polish Flag}

\paragraph{Problem} Given an array $A$ of $n$ balls, each of which is either red or white,
rearrange them in-place so that all the red balls are to the left.

\paragraph{Algorithm} Pretend that we have two elements, $A[-1]$ (this is the thing in the far left,
not the last element like in Python) as a dummy white ball, and $A[n]$ as a dummy red ball. We fix
two indices: $w=-1, r=n$. Notice we start right before the first elem and right after the last elem,
respectively. Why do we need these dummies though? \\
While $w \leq r$: swap $A[w] \iff A[r]$, increment $w$ until $A[w]$ is white, decrement $r$ until
$A[r]$ is red.

Invariants:
\begin{itemize}
    \item I1: $A[r]$ is red
    \item I2: nothing to the right of $A[r]$ is red (we want whites to the right... this is such
        a funny phrase it made me chuckle please don't cancel me I'm just trying to make my notes
        fun so I get an A)
    \item I3: $A[w]$ is white
    \item I4: nothing to the left of $A[w]$ is white (we want reds to the left)
\end{itemize}

\paragraph{Correctness}
\begin{proof}
    The first idea of the correctness proof is showing that our entire operation permutes the array.
    That is, our only whites are swaps $\implies$ permutation. Thus, we don't have to consider
    resizing or any alterations to the array, only proving the reds are to the left and the whites
    to the right (there it is again please laugh everyone!). \\
    \textbf{Initialization}: On initialization, I1, I2, I3, I4 are all vacuous, since there is 
    nothing to the right of $A[r]$ and nothing is to the left of $A[w]$. \\
    \textbf{Preserved}: We perform a swap, so I1, I3 are temporarily broken, but we fix this before
    the end of this iteration so I1 and I3 remain true. Our swap before incrementing, as well as our
    loop guard $w \leq r$ ensure I2, I4 remain true. \\
    \textbf{Correctness}: Once we terminate, we know our guard is false, that is $w > r$.
    At that point, I1, I2, I3, I4 are all satisfied, so we've reached our goal. \\
    $\mathbb{Q}$: At the end, how are $w$ and $r$ related? $\implies$ $w=r+1$. This implies that our
    algorithm \textbf{always} terminates! Note that proving this algorithm terminates is quite
    important and a very nontrivial step.
\end{proof}

\paragraph{Runtime}
\begin{proof}
    $w$ is incremented by $+1$ and $r$ is incremented by $-1$ for at most $n$ steps, so we have upper
    bound $O(n)$. As for our lower bound: $r-w=n+1$ at the start, and at the end $w=r+1 \implies 
    \Omega(n)$ lower bound runtime. Thus we have a tight bound of $\Theta(n)$.
\end{proof}

\subsection{Reductions}

\begin{definitionbox}{Reduction}
    A reduction from problem $P$ to problem $Q$ is an algorithm that solves $P$ using (any algorithm
    for) $Q$. That is, we can use $Q$ to solve $P$. We express this as $P \leq Q$ "$P$ reduces to
    $Q$. \\
    A nice way to think about this is taking an input $x$, (implicitly) transforming its input to 
    $Q$, then transforming the output to $P(x)$. \\
    Note that neither problem need be solvable!
\end{definitionbox}

\begin{itemize}
    \item Squaring $\leq$ multiply: $x^2 = x \times x$
    \item And the other way around: multiply $\leq$ square: $(x+y)^2-(x-y)^2$
    \item Median $\leq$ Sorting: sorted ($A$)$[\lfloor|A| / 2\rfloor]$
\end{itemize}

\subsubsection{Types of Reduction: Turing/Oracle and Many-one}

\begin{definitionbox}{Turing Reduction}
    We denote this as $\leq_T$. We can use $Q$ any number of times. Formally, a Turing Reduction from
    a problem $P$ to problem $Q$ is an algorithm $A$ that solves $P$ while using a subroutine solving
    $Q$ as an atomic operation.
\end{definitionbox}

\begin{definitionbox}{Many-one/Mapping Reduction}
    We denote this as $\leq_m$. A many-one reduction from problem $P$ to problem $Q$ is an algorithm 
    (or function) $f$ with the property that $P = Q \circ f$ or $P = Q(f(x))$ for input $x$. A many-
    one reduction is a Turing reduction with the restriction that the subroutine for $Q$ can be used
    only once, at the ennd, and the answer \textbf{must be returned without modification}. \\
    In other words, take an input $x$ for $P$, transform it using $f$ into an input $y=f(x)$ for $Q$,
    and then compute $Q(y)$, which should give the same answer as $P(x)$. A Many-One Reduction is
    necessarily a Turing Reduction.
\end{definitionbox}

\begin{itemize}
    \item Square $\leq$ mult: This is a turing reduction, since we use mult at least once. Since we
        only use mult once in $f(x)=x \times x$ this is also a many-one reduction.
    \item Mult $\leq$ square: We use square twice, so this is a turing reduction, but not a many-one
        reduction.
    \item Median $\leq$ Sort: It is a turing reduction, since we use it once, but it's not a many-one
        reduction for an odd reason: \textbf{the output type is different}!
\end{itemize}

\subsubsection{Median Reduces to Selection (Many-One)}
Assume for the input $A, |A|$ odd and all elements distinct. \\
A first instinct would be to recurse, perhaps dividing the array into three and recursing on the
middle partition, but that actually doesn't work, because the median of our subarray is not
necessarily the median of our larger array. Instead, let's reduce this to Selection.
\paragraph{Selection Algorithm}
\begin{itemize}
    \item input: array $A$ with distinct elems and an integer $k:0\leq k < |A|$.
    \item output: $k$th smallest elem of $A$.
\end{itemize}
For our selection algorithm, let's try the following. We guess an element $P \in A$ and check if that
is our $k$th smallest by colouring all smaller elements red and running Polish Flag. When Polish
terminates, we have 3 cases:
\begin{itemize}
    \item $w=k$: We have found our median, return $P$.
    \item $w>k$: Recurse on the left (red) side: SELECT($A[:w],k$). We know our $k$th smallest is
        in the smaller half of our array, so we can discard the upper half.
    \item $w<k$: SELECT($A[w+1:],k-w$). We know our $k$th smallest element is in the upper half of
        our array, so we can discard the lower half and increment $k$ by $w$ to adjust for our new
        subarray.
\end{itemize}
\paragraph{Now, how do we choose $P$?} A fixed (determinisitic) fast method for choosing $P$ implies
lower bound $\Omega(n^2)$, since the fixed algorithm may have some worsr case input for which it is 
ill-suited. Instead, let's choose \textbf{randomly}!

\subsection{Randomisation}

\subsubsection{Quick Select}
\begin{itemize}
    \item Base case: if $|A|$ small, just sort and return the index
    \item Choose random $i \in (0,|A|)$; set $P=A[i]$
    \item Colour the smaller elements red; run Polish Flag
    \item If $w=k$: output $P$
    \item If $w>k$: QUICKSELECT($A[:w], k$)
    \item If $w<k$: QUICKSELECT($A[w+1:], w-k$)
\end{itemize}

\paragraph{Why can we only claim partial correctness?} Exercise: Modify the algorithm slightly to
make it correct (What are we talking about here?)

\begin{definitionbox}{Expected Runtime}
    The expected runtime of a randomized algorithm is the maximum (over all inputs $x$ of length $n$)
    of the expected (over the randomness of our algorithm) number of steps that our algorithm
    executes for input $x$. Caution: We're looking at the worst-case input here, not a random input!
    Ensure to \textbf{quantify} the input first and then \textbf{take the expectation of the
    algorithm's runtime}. \\
    Note that the expected runtime is not the same as the average case runtime.
\end{definitionbox}

Let's prove the expected runtime of Quick select.

\begin{proof}
    Let pivot be the $i$th smallest element in $\{0,...,n-1\}$. Our expected runtime is the maximum
    number of steps (over $k$) that our algorithm executes. Now, when we choose a $k$, we have two
    cases: either our $k$ too large and we recurse on the ones before $k$, or our $k$ too small and
    we recurise on the ones after $k$. Note also that the probability of choosing this $k$ for which
    we have this maximum number of steps is $\frac{1}{n}$. Asymptotically, these sums are identical,
    so we can do the following
    \begin{align*}
        T(n) &\leq \max_k \Bigg[ \frac{1}{n}\sum_{i=0}^{k-1} T(n-i) 
            + \frac{1}{n}\sum_{i=k+1}^{n-1} T(i) + \Theta(n) \Bigg] \\
             &\leq \frac{2}{n}\Bigg( \sum_{i=n/2+1}^{n} T(i) \Bigg) + \Theta(n) \\
    \end{align*}
    Now, consider the pivot position (after Polish flag) landing in one of three regions:
    \begin{itemize}
        \item 0 to $n/4$: probability $1/4$
        \item $n/4+1$ to $3n/4$: probability $1/2$
        \item $3n/4+1$ to $n-1$: probability $1/4$
    \end{itemize}
    The first and last regions are pretty bad regions for our pivot to fall, since if we get unlucky
    we then have to consider at most $~n$ elements on our next recursive step. The middle region is
    pretty good, since if we get unlucky we only have to recurse on $3n/4$ elements on our next step.
    So, (considering our next recursive step) with probability $1/2$ we're in the $T(n)$ region and
    with probability $1/2$ we're in the $T(3n/4)$ region.
    \begin{align*}
        T(n) &\leq \frac{1}{2}T(n) + \frac{1}{2}T(3n/4) + \Theta(n) \\
             &\leq T(3n/4) + \Theta(n)
    \end{align*}


\end{proof}

\section{Lecture 5: Sequences}

\begin{definitionbox}{Data Structure}
    A data structure is a data store, plus a means of accessing and manipulating the stored data.
    Algorithms rely on data structures to compute with large amopunts of data and data structures
    are built on algorithms that tell them how to manipulate the data they store.
\end{definitionbox}

\begin{definitionbox}{Interface}
    An interface is a mathematical sepcification for a data structure. It states the type of
    object(s) represented and the required behaviour of supported operations. An implementation of
    an interface describes how to represent the data in memory, together with algorithms to perform
    each of the required operations. Note that an interface may have several different
    implementations - \textbf{they are not the same thing}!
\end{definitionbox}

\subsection{Stack Interface and its Implementations}
A stack uses a \textbf{last in first out} access pattern via
\begin{itemize}
    \item \texttt{push($x$)}: stores object $x$ on the top of the stack
    \item \texttt{peek()}: returns the object $x$ currently at the top of the stack
    \item \texttt{pop()}: returns and deletes the object $x$ currently at the top of the stack
    \item \texttt{is-empty()}: returns \textbf{TRUE} iff the stack has no objects
\end{itemize}

\subsubsection{Linked List}
\begin{itemize}
    \item Pointer-based data structure
    \item Each item stored in a \texttt{node} which contains a pointer to the next node below it.
        Each \texttt{node} has two fields: \texttt{node.item} and \texttt{node.next}
    \item Maintains a pointer to the top node \texttt{head} in the stack
    \item Insertion and deletion from the top takes $\Theta(1)$ time since we can just re-link
        pointers
    \item Can peek and check emptiness by reading \texttt{head} $\implies \Theta(1)$ time
\end{itemize}

\subsubsection{Static Array}
\begin{itemize}
    \item Arrays are contiguous chunks of memory cells.
    \item We maintain a pointer \texttt{top} (this isn't REALLY a pointer, but it's functionally one)
        pointing to the first empty cell in the array. Initialize \texttt{top} to 0.
    \item \texttt{push(c)}: we set \texttt{A[top]=c} and increment \texttt{top}.
    \item \texttt{pop()}: we decrement \texttt{top} and then return \texttt{A[top]}
    \item \texttt{peek()}: same as \texttt{pop()} then increment \texttt{top} to maintain our array
        state
    \item \texttt{is-empty()}: return \textbf{TRUE} iff \texttt{top} greater than 0
\end{itemize}

\subsubsection{Dynamic Array}
\begin{itemize}
    \item $n$ is the number of items stored and $m \geq n$ is the capacity. A resize that copies $n$
        items is linear, one that doesn't resize is constant.
\end{itemize}

\subsection{Amortized Analysis}
We consider the worst-case runtime of a sequence of operations. Now, this is a pretty nontrivial
topic, but I've found that the best explanation for me is the following amortized analysis of
appending (\texttt{push($x$)}) a sequence of values to a dynamic array.


\paragraph{\href{https://courses.cs.duke.edu/spring19/compsci330/lecture21scribe.pdf}{Proving Amortized Constant Time for insertion into dynamic array sequence:}}
Suppose we double the size of the array when full.  That is, we reallocate the entire array in memory when $i=2^k+1$. Let's analyze the total cost of $n$ insertions, then divide this total cost by $n$.

\begin{proof}
    \textbf{What is the cost of the i-th operation?} Case 1: If we do not need to increase the size 
    of the array, algorithm takes $O(1)$ unit of time. Case 2: If we need to increase the size of 
    the array, then at this step we must:
    \begin{enumerate}
        \item allocate memory for the new array (constant)
        \item copy over the $2^k$ elements from the old array to the new array ($2^k$ running time)
        \item free the memory storing the old array (constant)
        \item append the new element (constant)
    \end{enumerate}
    so this algorithm takes $2^k+1$ running time.
    Let $t(i):=$ the amount of running time.
    \begin{align*}
    t(i)= \begin{cases}
        2^k+1, \text{ if } i=2^k+1 \\
        1 \text{ otherwise}
    \end{cases}
    \end{align*}
    So we just sum over $n$.
    \begin{align*}
        \sum_{i=1}^{n} t(i) &= n + \sum_{{k=0}}^{\lfloor\log_2(n-1)\rfloor}2^k \\
        &= n+ \frac{2^{\lfloor\log_2(n-1)\rfloor}-1}{2-1} \\
        &\leq n+2n \\
        &\leq 3n
    \end{align*}
    So the total cost of $n$ insertions is at most $3n$, which has amortized cost at most $3$. 
    Thus, the amortized cost is $O(1)$.
\end{proof}

\subsection{Memory Overhead}

\begin{definitionbox}{Memory Overhead}
    Memory overhead is the memory used in addition to our memory allocated for our data, as a
    function of the number of elements stored. Dynamic arrays have an overhead of $\Theta(n)$ since
    we may allocate space $2n$ when storing $n$ elements (right after doubling). Linked list also
    have $\Theta(n)$ memory overhead, since in addition to the data stored in each node, we need to
    store a $\Theta(1)$ size memory address pointing to the next node in our list.
\end{definitionbox}

There's also the \texttt{SEQUENCE} interface, but I can't be bothered to write it here. You can
implement all the previous implementations within the sequence interface, with some
runtime advantages/tradeoffs and available constant time operations.


\section{Lecture 6: Sets and Hashes}
\begin{itemize}
    \item \texttt{Build()} (Initialize Container)
    \item \texttt{size()} (Container)
    \item \texttt{Find(e)} (Static)
    \item \texttt{Insert(x) - replace if already exists} (Dynamic)
    \item \texttt{Delete(k)} (Order)
\end{itemize}

\subsection{Comparison Model of Computation}
We associate a key with each object/item stored and determine equivalence by objects' keys. \\
One way to ensure constant time lookup is to use a \textbf{direct access array}. Say we have an 
item with key $k$, our item is stored in the $k-th$ position in the sequence. Insertion and deletion
also take constant time. However, we don't know how high our indexing goes. For example, if given 
9 digit integer keys, we'd need to allocate \textbf{$10^9$} slots in memory to store the items. 
Moreover, the size of our largest key $u$ is bounded by the word size $w$: $u < 2^w$ (this is 
strictly less than, because we're assuming $0 \leq k \leq 2^w-1$).

\paragraph{What if we surpass this bound?} We map a large space of $u$ keys down to a range 
$m=\Theta(n)$. In other words, we define a map:

\begin{align*}
    h: \{0,...,u-1\} \rightarrow \{0,...,n-1\}
\end{align*}

that maps the big space of $u$ things to a small space of $\Theta(n)$ things. By the pigeonhole 
principle, some of our items must collide after mapping. We have two options to deal with 
collisions:
\begin{enumerate}
    \item \textbf{Open addressing:} over allocate our $\Theta(n)$ array and find another place to 
        store it in there. (unimportant)
    \item At that key, store a pointer to another data structure (like a sequence), where we can 
        store a collection of items with collisions. That is, if our hashing function assigns $x$
        items the same key $k$, then at the corresponding position in our smaller array that we 
        map to, we can point to another data structure storing the $x$ items.
\end{enumerate}
Recall that by the pigeonhole principle, collisions are inevitable with any hash function and for 
$u$ very large, there is a possibility that all of our items collide. For now, we consider the 
easiest way to get from a large space of keys to a smaller one.

\subsubsection{Division Method}

\begin{align*}
    h(k)=k \text{mod(m)}
\end{align*}

This way, we take a key in the large space, and \textbf{mod} it so that it wraps around $m$ spaces.
In general, this distributes collisions uniformly given our keys are selected nicely. However, there
still exists some collection of inputs that map everything to one place. In any case, for any fixed
hash function, we can end up having all items collide (in the deterministic case). Instead, we can 
take a non-deterministic approach, where we choose a random hash function for any collection of 
keys input by the user.

\subsubsection{Universal Hash Function}
So, we define a family of hash functions and randomly choose from this collection. Start with a 
fixed hash function as shown below.

\begin{align*}
    h_{ab}=(((ak+b)\text{ mod p) mod m}) 
\end{align*}

Now, generalize to a family:

\begin{align*}
    H(p,m)=\{h_{ab}(k) | a,b \in \{0,...,p-1\} \text{ and } a \neq 0 \}
\end{align*}

We have a hash family parameterized by the length of our smaller array $m$ and fix some random 
large prime $p$ bigger than $u$. When we instantiate the hash table, we choose a random $a$ and 
$b$ to generate our function. By doing this \textbf{non-deterministic} jumbling, it is very 
unlikely for any given collection of keys to have too frequent collisions. Using this 
\textbf{universal hash family}, for any two keys, the probability that they collide is less than 
$1/m$. \\
Define random variable $X_{ij}$ over choice of $h \in H$, which means $X_{ij}=1$ if 
$h(k_i)=h(k_j)$, $0$ otherwise. \\
Size of chain at $h(k_i)$ is just the sum:

\begin{align*}
    X_i&=\sum_{j=0}^{n-1}X_{ij} \\
    \mathbb{E}_{h \in H} [X_i]&=\mathbb{E}\Bigg[\sum_{j}X_{ij}\Bigg] \\
    &=\sum_{j\neq i}\mathbb{E}[X_{ij}]+1 \\
    &= \sum_{j\neq i}^{n-1}(\frac{1}{m})+1 \\
    &= \frac{n-1}{m} + 1
\end{align*}

\section{}











\newpage

\section{Key Formulae and Theorems... Thanks ChatGPT}

\subsection{Theorems to Know}

\subsubsection{Master's Theorem}
For a recurrence of the form
\[
T(n) = aT\left(\frac{n}{b}\right) + f(n),
\qquad a \geq 1,\ b > 1,
\]
let
\[
c = \log_b a.
\]

Then:
\[
T(n) =
\begin{cases}
\Theta(n^c), & f(n) = O(n^{c-\epsilon}) \text{ for some } \epsilon > 0, \\[4pt]
\Theta(n^c \log n), & f(n) = \Theta(n^c), \\[4pt]
\Theta(f(n)), & f(n) = \Omega(n^{c+\epsilon}) \text{ for some } \epsilon > 0
\text{ and } af(n/b) \leq cf(n) \text{ for some } c<1.
\end{cases}
\]

\subsubsection{Fermat's Little Theorem}

If $p$ is prime and $a\not\equiv 0\pmod p$, then
\[
a^{p-1}\equiv 1\pmod p.
\]

Equivalently, for any integer $a$,
\[
a^p\equiv a\pmod p.
\]

\subsubsection{Bezout's Identity}

For integers $a,b$, there exist integers $x,y$ such that
\[
ax+by=\gcd(a,b).
\]

In particular,
\[
\gcd(a,b)=1
\iff
\exists x,y\in\mathbb Z
\text{ such that } ax+by=1.
\]

\subsubsection{Euclidean Algorithm}

For integers $a,b$ with $b>0$,
\[
\gcd(a,b)=\gcd(b,a\bmod b).
\]

Repeatedly applying this identity gives the Euclidean algorithm.

Its running time is
\[
O(\log\min(a,b)).
\]

\subsubsection{Extended Euclidean Algorithm}

The extended Euclidean algorithm computes integers $x,y$ satisfying
\[
ax+by=\gcd(a,b).
\]

If
\[
\gcd(a,m)=1,
\]
then the coefficient $x$ gives the modular inverse:
\[
ax\equiv 1\pmod m.
\]

\subsubsection{Euler's Theorem}

If
\[
\gcd(a,n)=1,
\]
then
\[
a^{\varphi(n)}\equiv 1\pmod n,
\]
where $\varphi(n)$ is Euler's totient function.

Fermat's Little Theorem is the special case
\[
n=p
\]
for prime $p$, since
\[
\varphi(p)=p-1.
\]

\subsubsection{Chinese Remainder Theorem}

If
\[
\gcd(n_i,n_j)=1
\qquad\text{for all }i\neq j,
\]
then the system
\[
x\equiv a_1\pmod{n_1},
\quad
x\equiv a_2\pmod{n_2},
\quad\ldots,\quad
x\equiv a_k\pmod{n_k}
\]
has a unique solution modulo
\[
N=\prod_{i=1}^k n_i.
\]

\subsubsection{Binomial Theorem}

For any integer $n\geq 0$,
\[
(x+y)^n
=
\sum_{k=0}^{n}
\binom{n}{k}x^{n-k}y^k.
\]

The binomial coefficient is
\[
\binom{n}{k}
=
\frac{n!}{k!(n-k)!}.
\]

\subsubsection{Pascal's Identity}

For $0<k<n$,
\[
\binom{n}{k}
=
\binom{n-1}{k-1}
+
\binom{n-1}{k}.
\]

Boundary cases:
\[
\binom{n}{0}=\binom{n}{n}=1.
\]

\subsubsection{Handshaking Lemma}

For any undirected graph $G=(V,E)$,
\[
\sum_{v\in V}\deg(v)=2|E|.
\]

Consequently, the number of vertices of odd degree is even.

\subsubsection{Pigeonhole Principle}

If $n+1$ objects are placed into $n$ boxes, then at least one box contains at least two objects.

More generally, if $N$ objects are placed into $k$ boxes, then some box contains at least
\[
\left\lceil\frac Nk\right\rceil
\]
objects.

\subsubsection{Inclusion--Exclusion Principle}

For two sets,
\[
|A\cup B|
=
|A|+|B|-|A\cap B|.
\]

For three sets,
\[
\begin{aligned}
|A\cup B\cup C|
={}&|A|+|B|+|C|\\
&-|A\cap B|-|A\cap C|-|B\cap C|\\
&+|A\cap B\cap C|.
\end{aligned}
\]

In general,
\[
\left|\bigcup_{i=1}^n A_i\right|
=
\sum_{\emptyset\neq S\subseteq[n]}
(-1)^{|S|+1}
\left|\bigcap_{i\in S}A_i\right|.
\]

\subsubsection{Cauchy--Schwarz Inequality}

For real numbers $a_1,\ldots,a_n$ and $b_1,\ldots,b_n$,
\[
\left(\sum_{i=1}^n a_i b_i\right)^2
\leq
\left(\sum_{i=1}^n a_i^2\right)
\left(\sum_{i=1}^n b_i^2\right).
\]

\subsubsection{Triangle Inequality}

For real numbers,
\[
|a+b|\leq |a|+|b|.
\]

More generally,
\[
\left|\sum_{i=1}^n a_i\right|
\leq
\sum_{i=1}^n |a_i|.
\]

\subsubsection{Law of Total Probability}

If $B_1,\ldots,B_n$ form a partition of the sample space, then
\[
P(A)
=
\sum_{i=1}^n P(A\mid B_i)P(B_i).
\]

\subsubsection{Bayes' Theorem}

For events $A$ and $B$ with $P(B)>0$,
\[
P(A\mid B)
=
\frac{P(B\mid A)P(A)}{P(B)}.
\]

More generally, if $B_1,\ldots,B_n$ form a partition,
\[
P(B_j\mid A)
=
\frac{P(A\mid B_j)P(B_j)}
{\sum_{i=1}^nP(A\mid B_i)P(B_i)}.
\]

\subsubsection{Linearity of Expectation}

For random variables $X_1,\ldots,X_n$,
\[
\mathbb E\left[\sum_{i=1}^nX_i\right]
=
\sum_{i=1}^n\mathbb E[X_i].
\]

This holds whether or not the random variables are independent.

\subsubsection{Markov's Inequality}

For a nonnegative random variable $X$ and $a>0$,
\[
P(X\geq a)
\leq
\frac{\mathbb E[X]}{a}.
\]

\subsubsection{Chebyshev's Inequality}

For a random variable $X$ with finite mean $\mu$ and variance $\sigma^2$,
\[
P(|X-\mu|\geq t)
\leq
\frac{\sigma^2}{t^2}.
\]

Equivalently,
\[
P(|X-\mu|\geq k\sigma)
\leq
\frac1{k^2}.
\]

\subsubsection{Union Bound}

For any events $A_1,\ldots,A_n$,
\[
P\left(\bigcup_{i=1}^n A_i\right)
\leq
\sum_{i=1}^nP(A_i).
\]

\subsubsection{Geometric Distribution}

If independent trials have success probability $p$, then the probability that the first success occurs on trial $k$ is
\[
P(X=k)=(1-p)^{k-1}p.
\]

Its expectation is
\[
\mathbb E[X]=\frac1p.
\]

\subsubsection{Birthday Bound}

For $k$ uniformly random choices from a set of size $n$, the probability of at least one collision is approximately
\[
1-e^{-k(k-1)/(2n)}.
\]

In particular, collisions become likely when
\[
k=\Theta(\sqrt n).
\]

\subsection{Summation Formulae}
% =========================
% Summation Formulae - should probably partition this better
% =========================

\subsubsection{Basic Summations}

\[
\sum_{i=1}^{n} 1 = n
\]

\[
\sum_{i=1}^{n} i = \frac{n(n+1)}{2} = \Theta(n^2)
\]

\[
\sum_{i=1}^{n} i^2
= \frac{n(n+1)(2n+1)}{6}
= \Theta(n^3)
\]

\[
\sum_{i=1}^{n} i^3
= \left(\frac{n(n+1)}{2}\right)^2
= \Theta(n^4)
\]

\[
\sum_{i=1}^{n} i^k = \Theta(n^{k+1})
\qquad (k > -1)
\]

% =========================
% Arithmetic / Geometric
% =========================

\subsubsection{Arithmetic Series:}

\[
\sum_{i=0}^{n-1} (a+id)
=
\frac{n}{2}\left(2a+(n-1)d\right)
\]

\subsubsection{Geometric Series:}

\[
\sum_{i=0}^{n-1} r^i
=
\frac{1-r^n}{1-r}
\qquad (r\neq 1)
\]

\[
\sum_{i=0}^{\infty} r^i
=
\frac{1}{1-r}
\qquad (|r|<1)
\]

\[
\sum_{i=0}^{n} r^i = \Theta(r^n)
\qquad (r>1)
\]

% =========================
% Harmonic
% =========================

\subsubsection{Harmonic Series:}

\[
H_n
=
\sum_{i=1}^{n}\frac{1}{i}
=
\Theta(\log n)
\]

More precisely,

\[
H_n = \ln n + \gamma + O(1/n)
\]

\[
\sum_{i=1}^{n}\frac{1}{i^p}
=
\begin{cases}
\Theta(1), & p>1,\\
\Theta(\log n), & p=1,\\
\Theta(n^{1-p}), & 0\leq p<1.
\end{cases}
\]

% =========================
% Logarithmic
% =========================

\subsubsection{Logarithmic Summations:}

\[
\sum_{i=1}^{n}\log i
=
\log(n!)
=
\Theta(n\log n)
\]

\[
\sum_{i=1}^{n}\log(i)
=
\Theta(n\log n)
\]

\[
\sum_{i=1}^{n}\frac{\log i}{i}
=
\Theta((\log n)^2)
\]

% =========================
% Powers of Two
% =========================

\subsubsection{Powers of Two:}

\[
\sum_{i=0}^{\log_2 n} 2^i
=
2^{\log_2 n+1}-1
=
2n-1
=
\Theta(n)
\]

\[
\sum_{i=0}^{\log_2 n} i
=
\Theta((\log n)^2)
\]

\[
\sum_{i=0}^{\log_2 n} \frac{1}{2^i}
=
\Theta(1)
\]

\[
\sum_{i=0}^{\log_2 n} \frac{n}{2^i}
=
\Theta(n)
\]

% =========================
% Telescoping
% =========================

\subsubsection{Telescoping Sums:}

\[
\sum_{i=1}^{n}(a_i-a_{i+1})
=
a_1-a_{n+1}
\]

A common example:

\[
\sum_{i=1}^{n}\frac{1}{i(i+1)}
=
\sum_{i=1}^{n}
\left(\frac1i-\frac1{i+1}\right)
=
1-\frac1{n+1}
=
\Theta(1)
\]

% =========================
% Useful Asymptotic Rules
% =========================

\subsubsection{Useful Rules:}

\[
\sum_{i=1}^{n}\Theta(f(i))
=
\Theta\left(\sum_{i=1}^{n}f(i)\right)
\]

\[
\sum_{i=1}^{n}c\,f(i)
=
c\sum_{i=1}^{n}f(i)
\]

\[
\sum_{i=1}^{n}(f(i)+g(i))
=
\sum_{i=1}^{n}f(i)
+
\sum_{i=1}^{n}g(i)
\]

\[
\sum_{i=1}^{n} f(i)
\leq
n\max_{1\leq i\leq n} f(i)
\]

\[
\sum_{i=1}^{n} f(i)
\geq
n\min_{1\leq i\leq n} f(i)
\]

% =========================
% Integral Bounds
% =========================

\subsubsection{Integral Approximation:}

For nonnegative increasing $f$,

\[
\int_1^n f(x)\,dx
\leq
\sum_{i=1}^{n}f(i)
\leq
f(n)+\int_1^n f(x)\,dx
\]

For nonnegative decreasing $f$,

\[
\int_1^{n+1} f(x)\,dx
\leq
\sum_{i=1}^{n}f(i)
\leq
f(1)+\int_1^n f(x)\,dx.
\]

\end{document}
